\section{A Beautiful Day to Die}

\begin{quotex}
Even without being killed a man can experience death, he can conquer, he can realize the culmination characteristic of a ``super-life''. From a higher point of view, Paradise, the Kingdom of Heaven, Valhalla, the Island of the Heroes, etc., are only symbolic figurations forged for the masses, figurations that in reality designate transcendent states of consciousness, beyond life and death. The ancient Aryan tradition used the term \emph{jivan-mukti} to indicate such a realization while still in the mortal body.
\flright{\textsc{Julius Evola}, \emph{The Aryan Conception of Combat}}
\end{quotex}


A couple of weeks ago I had a dream in one scene of which I placed a book under my seat. When I reached down to pick it up, I noticed it had disappeared. Still in that dream state, I thought, ``This must be a dream, since in real life books don't disappear just like that.'' That is when I awoke. Similarly, the waking state of consciousness is itself revealed to be a dream from the higher state of Samadhi. Does anyone pay attention to the clues like the missing book to realize that?

If you believe your life to be a random accident, of no significance, there is nothing to notice. However, most men look for their purpose horizontally, as some goal to achieve within that dream. Have you noticed how your dreams are mostly incoherent, that there seems to be relationships from scene to scene, but turn out to be tenuous, that one's intention fails and morphs into another? Is that the description of a dream or life as you know it? If the latter, how do you know when it is time to wake up? When do you choose to wake up?

So, if you don't believe your life is random, and you see that your purpose is not found within life itself, and that your real life is your super-life, then how much effort are you putting into living your real life? Shouldn't you be obsessed with it?

Don't think you hide from it or evade this task, for you will reach that state at the moment of death. However, at that point, it will not happen consciously and freely, but rather will be determined and automatic. This man dies in the first or second state of consciousness. That is, in the first, he lives his life in a passive sense, thoughts flow through him. In the second, his mind is full of theories, opinions, philosophies that maintain him in the illusions of ordinary life. The jivan-mukti dies to eternal life, beyond life and death, in a transcendent state. The other man just dies in the state he found himself in.

\paragraph{The Last Thought}


That is why a man's final thoughts have been considered so important. Tibetan Buddhism has an elaborate ritual to guide a man toward his approaching death. The Greek and Roman Churches offer the sacrament of extreme unction to improve the quality of a man's final thought.

Observe your breath. Notice that every exhalation may be your last; you cannot be certain that you will inhale another. What are you thinking at that moment? Is it a thought you want to be remembered by? Or, better said, is it a thought you want to die remembering?

Some may think they are superior because they don't believe the sensual understanding of the post-mortem state in folk religion. However, Evola changes nothing in his esoteric interpretation. Even the popes have called Heaven a state of being, not a place, so the goal of a happy death remains.

\paragraph{The Samurai's Servant}


In Japanese legends, the last thoughts of a dying man were believed to have irresistible powers. Lafcadio Hearn tells the story of a samurai who condemned one of his slaves to death by beheading. The slave, convinced of the injustice of the sentence, bitterly told his master, at his execution, that he would take revenge. The samurai's family were terrified, because they understood the power of a dying man's final thoughts.

The clever samurai told the slave he didn't believe his threats unless the slave could prove it to him. The slave offered to prove it. So the samurai told him to bite the step as his head was rolling down; he would take that as proof of the slave's threat. The slave exulted, ``I will do so!'' The samurai then chops off the slave's head, which then latches onto the step with its teeth as it is rolling down.

The samurai's family and retinue were horrified to witness that proof of the slave's threat of vengeance. The samurai himself was unperturbed and gloated: ``The slave's final thought was to bite the step, which he did, so his earlier threat is ineffective!''

\url{https://www.youtube.com/watch?v=RnCmeNa5vXs}

\flrightit{Posted on 2013-01-30 by Cologero}

\begin{center}* * *\end{center}

\begin{footnotesize}
\begin{sffamily}
\texttt{A on 2013-01-31 at 00:38 said:}

Have you ever considered how would that samurai's final thought be after he has committed injustice and cunningness himself in his life.

\hfill

\texttt{Cologero on 2013-01-31 at 00:43 said:}

The samurai was a warrior, violent by nature, beyond good and evil. Does that count for anything?

\hfill

\texttt{Mihai on 2013-01-31 at 09:06 said:}

You know, this is a good exercise for when dealing with Thanthos' little brother- Hypnos. Before falling asleep -- to meditate on the possibility of never waking up again and focus all your attention on the name of God, with all the power of your being.

\hfill

\texttt{Janus on 2013-02-01 at 13:03 said:}

Well, the samurai as kshatriya would surely be required to conform that nature to the higher Self according to the rites of the warrior? He doesn't seem to be the quite so honourable samurai, although if we look at the Eddas then we have many similar stories of trickery. The tale nowhere says that the samurai is honourable or perfect, I suppose.

\hfill

\texttt{Cologero on 2013-02-01 at 13:15 said:}

I am glad this topic has come up. On Evolian grounds, how can you judge the Samurai? He is following the way of action that Evola claims is as valid a path to the ``higher self'' as the way of contemplation. In other words, action is its own judge. The Samurai is to act honorably to his peers and to be loyal to his superiors. However, he has no obligation to be honorable to the slave. Perhaps he should have followed ``noblesse oblige'' and acted magnanimously. So perhaps you could say the Samurai did not act nobly.

\hfill

\texttt{Janus on 2013-02-02 at 16:19 said:}

How would you say the Samurai has no duty to be honourable to the slave? Surely he must be honourable to everyone, whatever that may mean according to their relationship. In the case of the slave, that would mean dispensing true justice. If he did so and the slave was wrong, then magnanimity does come into it, I suppose. I'm not too sure how far the concept of noblesse oblige influenced the Samurai. As for Evola, this might be an example against his concept of contemplation and action being on equal planes. If the samurai's actions are not guided by the precepts of Bushido, founded on higher rites, principles and realities, then he is acting mindlessly (or perhaps better said, without Intellect), something which does not exactly communicate honour or discipline.

\hfill

\texttt{Cologero on 2013-02-02 at 18:02 said:}

Exactly, Janus, that is the point I didn't make so well. Does Evola's Nietzscheanism, action beyond good and evil, not subject to a higher law, and even embracing ``cruelty'', justify the Samurai? Evola is correct that there is a karma yoga of action, alongside jnana yoga, but is his understanding of it distorted? Next week I'll post something on Coomaraswamy's critique of Evola and perhaps we can revisit this topic with some additional conceptual tools.


\hfill

\end{sffamily}
\end{footnotesize}

